\documentclass[11pt]{article}\usepackage[margin=2.3cm]{geometry}\usepackage{enumitem}\setlength{\parindent}{0pt}\setlength{\parskip}{4pt}
\begin{document}
\section*{Response to reviewers --- POWTEC-D-26-03167 (revised version)}
We thank both reviewers. Re-running every analysis with a stricter protocol (fold-wise calibration of $\Pi_{VC}$, repeated geometry-grouped splits, ablation, LOGO/L2GO) showed that several claims of the first submission were not supported; we have withdrawn or qualified them and reframed the paper as a validation and ablation study. Numbers below refer to the revised manuscript.

\textbf{Correction made in this revision:} the feature vector contains 17 inputs (7 raw, 4 logarithms, 6 derived), not 14 as stated before; the exact list is in Sec.~3.1.

\subsection*{Reviewer 1}
\begin{enumerate}[leftmargin=*,label=(\arabic*)]
\item \textbf{Ablation.} Added (Sec.~4.2, Tables 5--6, Fig.~2): identical models with raw vs.\ 17 features, fixed and tuned hyperparameters, plus paired LOGO tests. Result: small consistent gains with tuned boosters (RMSE $-0.001$ to $-0.011$), none significant under LOGO, large gain only for the neural network. The claim of a methodological novelty has been removed; the paper is now framed as a validation study.
\item \textbf{Feature importance.} Agreed. The permutation-importance claim (40\%, 73\%) is deleted; the ablation replaces it.
\item \textbf{Empirical constants.} Agreed. $\Pi_{VC}$ is now explicitly a data-fitted group. Constants are refitted inside each training fold (Sec.~4.3); they are poorly identified, and the former values $a=0.70$, $b=0.30$, $c=2$ are not reproducible out-of-fold.
\item \textbf{Re scaling.} Eq.~(1) is corrected (missing $\tau_p$) and the $\sqrt{\mathrm{Re}}$ path-length argument is presented as an assumption. Fitting the exponent freely gives $+0.43$, not $-0.5$ (Sec.~4.3): the assumption is not supported and is reported as such.
\item \textbf{Weak generalization.} Accepted. We now quantify it: random split $R^2\approx0.95$, grouped splits and LOGO $\approx0.73$--0.75 (RMSE 0.17), with pooled metrics because per-geometry $R^2$ is unstable on small geometries (Sec.~4.1, 4.4). The abstract and conclusions state this directly.
\item \textbf{Mahalanobis criterion.} Re-tested on all 28 geometries: Spearman 0.13 ($p=0.50$), AUC 0.62 (95\% interval 0.33--0.89). It is not validated; the practitioner rule is withdrawn and the indicator is kept only as a descriptive distance (Sec.~4.6).
\item \textbf{``Factor 2'' widening.} Tested and not supported: required inflation is $\approx1.3$ and equal for flagged and unflagged geometries. Removed; replaced by out-of-geometry coverage results and a conformal rescaling (Sec.~4.7).
\item \textbf{Cut-diameter scaling.} Re-done per condition (73 curves) with geometry-clustered bootstrap and three specifications, including geometry fixed effects. Re and geometry are 98.7\% collinear; the exponent ($-0.46$ to $-0.55$) is reported as consistent with $-0.5$ but not separable from geometry effects, and no physical interpretation is claimed (Sec.~4.8).
\item \textbf{Pressure drop.} Agreed it does not support the method. The Schmid--Vogel data are not available in the supplementary files, so the pressure-drop section is removed.
\item \textbf{CatBoost vs.\ GBR.} Agreed. No algorithm is claimed to be best; all tree ensembles are equivalent within split-to-split spread (Sec.~4.1).
\end{enumerate}

\subsection*{Reviewer 2}
\begin{enumerate}[leftmargin=*,label=(\arabic*)]
\item \textbf{Dataset consistency.} Sec.~2 describes the construction (977 unique points after removing 3 duplicates; geometry and condition definitions), corrects the ranges in Table 1, and Appendix A characterizes each geometry. Source labels are not in the database, so inter-source bias cannot be modelled; this is stated as a limitation.
\item \textbf{Random vs.\ LOGO gap.} Analysed in Sec.~4.1 and 4.4 (Table 3): the gap comes mainly from point-wise leakage between points of the same curve; L2GO added.
\item \textbf{Training/hyperparameters.} Sec.~3.2 specifies two protocols (fixed identical, and tuned with identical budget and grouped inner CV).
\item \textbf{Re dependence.} Sec.~4.5 and Table 7: error by band, band composition, and leave-band-out tests. The applicability range is stated.
\item \textbf{References 3 and 4.} Full entries are now given for references 3--5 (to be checked by the authors against the publisher pages before submission).
\item \textbf{Consistency.} All numbers are regenerated by a single script; 0.956/0.947/0.955 and 0.993/0.995 no longer appear.
\item \textbf{GBR vs.\ framework.} The paper now proposes neither GBR nor a best model; the title and abstract are changed accordingly.
\item \textbf{Ablation.} See R1(1).
\item \textbf{$\Pi_{VC}$ determination and leakage.} Constants are fitted on the training data of each fold only (Sec.~3.1); the earlier global calibration is abandoned.
\item \textbf{Further geometry-level validation.} 20 repeated random 23/5 splits, LOGO, L2GO (60 pairs) and leave-Re-band-out (Sec.~3.3, 4).
\item \textbf{Mahalanobis evidence.} See R1(6): the evidence is insufficient and the claim is withdrawn.
\item \textbf{Additional literature.} Added in the Introduction and Discussion; the works are now discussed (hydrocyclone optimization/control, CNN--LSTM surrogates, hybrid models) and contrasted with the cross-geometry transfer addressed here.
\end{enumerate}
\end{document}
